\documentclass[12pt,a4paper]{article}

\usepackage{fontawesome5}
\providecommand{\faPencilAlt}{\faIcon{pencil-alt}}

\usepackage{styles/couleurs}
\usepackage{styles/mise_en_page}
\usepackage{styles/boites}
\usepackage{styles/maths}
\usepackage{styles/figures}

\usepackage{enumitem}
\usepackage{array}
\usepackage{tabularx}
\usepackage{amsmath}
\usepackage{amssymb}
\usepackage{tikz}
\usepackage{xparse}

\usetikzlibrary{arrows.meta,calc}

\setlist[itemize]{leftmargin=1.4em,itemsep=0.15em,topsep=0.2em}
\setlist[enumerate]{leftmargin=*,itemsep=0.25em,topsep=0.25em}

\newcommand{\lignepointillee}{%
  \noindent\dotfill\par
  \vspace{0.45em}
}

\newcommand{\lignesreponse}[1]{%
  \begin{tcolorbox}[
    enhanced,
    colback=white,
    colframe=NeutreC!40,
    boxrule=0.5pt,
    arc=1.5mm,
    left=3mm,
    right=3mm,
    top=2mm,
    bottom=1mm,
    before skip=2pt,
    after skip=2pt
  ]
    \foreach \i in {1,...,#1}{%
      \lignepointillee
    }
  \end{tcolorbox}
}

\newcommand{\titreseq}[1]{%
\begin{center}
  {\sffamily\large\bfseries\color{SavoirC}#1}
\end{center}
}

\begin{document}

\titreseq{Activité 6 -- Enchaînements avec les quatre opérations}

\section*{Partie A -- Trois copies}

On veut calculer :

\[
A=(-24)\div(-6)+5\times(-2).
\]

\begin{minipage}[t]{0.31\textwidth}
\textbf{Copie 1}
\[
A=4+(-10)=-6.
\]
\end{minipage}
\hfill
\begin{minipage}[t]{0.31\textwidth}
\textbf{Copie 2}
\[
A=(-24)\div(-1)\times(-2)=-48.
\]
\end{minipage}
\hfill
\begin{minipage}[t]{0.31\textwidth}
\textbf{Copie 3}
\[
A=(-4)+(-10)=-14.
\]
\end{minipage}

\begin{enumerate}
  \item Quelle copie est correcte ?
  \item Explique l'erreur des deux autres copies.
\end{enumerate}

\lignesreponse{6}

\section*{Partie B -- La méthode}

Complète.

\begin{tcolorbox}[colback=SavoirC!5,colframe=SavoirC!60,boxrule=0.7pt,arc=2mm]
Dans un enchaînement, on calcule d'abord les \dotfill

\vspace{0.4cm}

Puis les \dotfill{} et les \dotfill{} de gauche à droite.

\vspace{0.4cm}

Enfin les \dotfill{} et les \dotfill{} de gauche à droite.
\end{tcolorbox}

\section*{Partie C -- Corriger une erreur}

Un élève écrit :

\[
18-(-3)\times4=(18-(-3))\times4=21\times4=84.
\]

Explique l'erreur et corrige.

\lignesreponse{5}


\section*{Partie D -- Calculer en rédigeant}
\[
B=5+(-3)\times4
\]

\[
C=7-(-4)\times(-3)
\]

\[
D=\bigl(8-3\times(-2)\bigr)\div(-7)
\]

\[
E=(-3)\times(4-7)-18\div(-6)
\]

\noindent
\begin{minipage}[t]{0.49\linewidth}
\lignesreponse{26}
\end{minipage}\hfill
\begin{minipage}[t]{0.49\linewidth}
\lignesreponse{26}
\end{minipage}

\end{document}
